\chapter{ДИЗАЙН ЭКСПЕРИМЕНТОВ}

\section{Постановка экспериментальной задачи}

В первой главе была сформулирована гипотеза о том, что эффективность ClawdLab в задачах научного анализа и проработки гипотез может быть повышена за счет двух доработок: структурированной памяти исследования и управляемой циклической модели выполнения задач. Во второй главе эта идея переводится в экспериментальную постановку: определяется, на каких типах задач она проверяется, какие ограничения принимаются, какие конфигурации системы сравниваются и по каким критериям принимается решение о результате эксперимента.

\subsection{Типы научных задач, поддерживаемые ClawdLab}

ClawdLab ориентирован на поддержку исследовательского процесса, в котором несколько агентов и человек совместно формируют, обсуждают, критикуют и проверяют научные идеи [1]. Система представляет собой лабораторную среду, где агенты с различными ролями помогают организовать работу с исследовательскими идеями, источниками, аргументами, проверками и итоговыми артефактами. Эта система не является автономным ученым, выполняющим любые виды экспериментов без участия человека.

С учетом этого в рамках данной работы рассматриваются не все возможные научные задачи, а те типы задач, которые соответствуют назначению ClawdLab и могут быть воспроизведены в эксперименте:

- анализ научной информации и выделение исследовательского разрыва,
- сравнение существующих подходов по заданным критериям,
- генерация и уточнение исследовательских гипотез,
- критика предложенной гипотезы или плана проверки,
- подготовка структурированного вывода на основе набора источников и промежуточных артефактов.

Такие задачи соответствуют ранним и средним этапам научного процесса: обзору предметной области, постановке проблемы, формулированию гипотез и предварительной проверке их обоснованности. Они также близки к сценариям, рассматриваемым в Agent Laboratory, где исследовательский процесс включает обзор литературы, экспериментирование и подготовку отчета [2].

При этом объектом экспериментального сравнения в данной работе являются не все перечисленные операции по отдельности, а два типа задач ClawdLab: analysis и deep research. Остальные типы, такие как literature review, critique и synthesis, рассматриваются как стадии или вспомогательные действия для выполнения этих задач. Это позволяет сохранить связь с архитектурой ClawdLab, но не расширять эксперимент за пределы проверяемых доработок.

\subsection{Границы эксперимента}

Эксперимент не направлен на проверку способности системы выполнять полный автономный цикл научного открытия во всех дисциплинах. В частности, в рамках работы не рассматриваются задачи, требующие физического эксперимента, сложной лабораторной инфраструктуры, длительного обучения моделей или доступа к закрытым данным.

Также в экспериментальную архитектуру не включается динамический выбор языковых моделей. Это связано с тем, что ClawdLab предполагает подключение внешних агентов с уже заданными моделями и исполнительными контурами. Платформа может регистрировать агента, задавать его роль и координировать взаимодействие, но не управляет напрямую выбором модели внутри внешнего агента. Поэтому выбор моделей фиксируется на уровне экспериментальных условий, а не рассматривается как самостоятельная доработка.

Фокус эксперимента ограничен двумя архитектурными свойствами:

- наличие или отсутствие структурированной памяти исследования,
- наличие или отсутствие циклической модели выполнения задач.

Такое ограничение позволяет проверить именно те доработки, которые могут быть встроены в ClawdLab как слой управления исследовательским состоянием и артефактами, не нарушая идею системы как платформы для координации внешних агентов.

\subsection{Исследовательский вопрос и ожидаемые эффекты}

Основной исследовательский вопрос формулируется следующим образом:

В какой степени добавление структурированной памяти исследования и управляемой циклической модели выполнения задач влияет на качество, согласованность и воспроизводимость результатов ClawdLab в задачах научного анализа и проработки гипотез?

Из этого вопроса следуют два ожидаемых эффекта, которые далее проверяются через ablation-сравнение. Первый ожидаемый эффект связан со структурированной памятью: явное связывание гипотез, источников, попыток, критики и выводов должно уменьшить потерю контекста и повысить согласованность итогового результата с промежуточными артефактами. Второй ожидаемый эффект связан с циклической моделью: управляемые попытки выполнения должны помочь системе использовать критику и исправлять неполные, противоречивые или недостаточно обоснованные результаты.

Формальная нулевая и альтернативная гипотезы задаются далее в разделе методики оценки. Такое разделение позволяет не смешивать содержательное описание ожидаемых эффектов с критерием статистической или качественной валидации.

Источники:

1. From Agent-Only Social Networks to Autonomous Scientific Research: Lessons from OpenClaw and Moltbook, and the Architecture of ClawdLab and Beach.Science - [https://arxiv.org/abs/2602.19810](https://arxiv.org/abs/2602.19810)
2. Agent Laboratory: Using LLM Agents as Research Assistants - [https://arxiv.org/abs/2501.04227](https://arxiv.org/abs/2501.04227)


\section{Обзор инструментов, технологий и методов для экспериментальной реализации}

В данном разделе рассматриваются не все технологии, связанные с LLM-агентами, а только те подходы, которые могут быть использованы для экспериментальной проверки двух выбранных доработок: структурированной памяти и циклического выполнения задач. В отличие от первой главы, где основное внимание уделялось related work и архитектурным направлениям, здесь анализируются практические варианты реализации, пригодные для прототипа.

\subsection{Требования к выбираемым технологиям}

Технологические решения для эксперимента должны соответствовать нескольким требованиям. Во-первых, они должны быть совместимы с текущей логикой ClawdLab, где агенты являются внешними исполнителями, а сама платформа отвечает за координацию, хранение артефактов и управление исследовательским процессом. Во-вторых, они должны позволять сравнивать базовую и улучшенную конфигурации в одинаковых условиях. В-третьих, они не должны чрезмерно усложнять прототип, поскольку цель эксперимента состоит не в построении полной production-архитектуры, а в проверке сформулированных ожидаемых эффектов.

Из этого следует, что предпочтение отдается эволюционным доработкам, которые могут быть добавлены поверх существующих сущностей: задач, обсуждений, артефактов, документов и журналов активности. Более сложные подходы, например полная графовая память с формальной логикой пересмотра убеждений, рассматриваются как ориентир для развития, но не обязательно как минимальная реализация эксперимента.

\subsection{Технологии структурированной памяти}

Существует несколько подходов к организации памяти: событийная память, векторная память, иерархическая память, графовая память и артефактно-ориентированное состояние. Для экспериментальной реализации наиболее важным является не само наличие хранилища, а возможность связать между собой элементы научного процесса: гипотезу, задачу, попытку выполнения, результат, критику и решение о дальнейших действиях.

Самый простой вариант памяти в агентных системах - журнал событий или история сообщений. Такой подход технически реализуется как последовательность логов: агент получает предыдущие сообщения, результаты вызовов инструментов и текстовые артефакты. Его преимущество состоит в простоте реализации и хорошей совместимости с существующими системами. Однако такая память плохо масштабируется: по мере роста числа действий агенту становится трудно отделять важные выводы от служебных деталей, а связи между гипотезами, экспериментами и выводами остаются неявными.

Более развитый вариант - извлекаемая память на основе векторного поиска (retrieval-based memory). В этом случае отдельные фрагменты текста, документов или результатов преобразуются в эмбеддинги и сохраняются в векторном хранилище. При выполнении новой задачи система извлекает релевантные фрагменты по семантической близости. Такой подход полезен для поиска похожих документов или прошлых обсуждений, но он не всегда показывает, почему именно один вывод связан с другим. Для научной задачи этого недостаточно: важно хранить не только близость текстов, но и тип отношения между ними, например "поддерживает", "опровергает", "уточняет" или "основано на эксперименте".

Графовый подход представлен в работе Young Bin Park "Graph-Native Cognitive Memory for AI Agents" [1]. Автор предлагает архитектуру Kumiho, в которой память строится как графовая когнитивная память (graph-native cognitive memory). В такой архитектуре узлы графа могут соответствовать утверждениям, наблюдениям, источникам, гипотезам или версиям артефактов, а ребра фиксируют зависимости: поддержку, противоречие, происхождение, уточнение или пересмотр. Важной особенностью является версионирование памяти: вместо перезаписи старого состояния создаются неизменяемые ревизии (immutable revisions), но актуальные указатели на них могут меняться. Также используется двойное хранилище: рабочая память для быстрых операций и долгосрочный граф для устойчивого хранения связей.

Для научных задач графовая память является полезной, потому что она позволяет фиксировать происхождение вывода и отношения между доказательствами. Например, гипотеза может быть связана с набором источников, экспериментальными результатами и критическими замечаниями. Если новая информация опровергает часть гипотезы, система может не просто добавить еще один текстовый фрагмент, а явно зафиксировать отношение противоречия и необходимость пересмотра. Однако полная реализация такого подхода требует отдельной семантической модели, правил версионирования, поддержки графового хранилища и механизма пересмотра убеждений. Для первого прототипа ClawdLab это решение является полезным ориентиром, но слишком сложным как минимальная доработка.

Более практичным для текущей работы является иерархический подход к памяти. В статье Xinyu Zhu и соавт. "Toward Ultra-Long-Horizon Agentic Science" предложена технология Hierarchical Cognitive Caching (HCC), используемая в системе ML-Master 2.0 [2]. Авторы рассматривают управление контекстом как процесс когнитивного накопления (cognitive accumulation), его задача состоит в том, чтобы отделить краткосрочные детали выполнения от долгосрочной исследовательской стратегии.

Технически HCC организует память в несколько уровней. На нижнем уровне сохраняются временные следы выполнения: действия агента, ошибки, результаты запусков, параметры и наблюдения. Эти данные слишком подробны, чтобы постоянно передавать их в контекст модели. Поэтому система выполняет дистилляцию: из временных следов выделяются более устойчивые эпизоды, выводы и уроки. На верхних уровнях хранятся уже не все детали, а сжатые знания, которые могут использоваться при планировании следующих шагов. Таким образом, HCC превращает поток выполнения в многоуровневую память: рабочую, эпизодическую и долгосрочную.

Применительно к ClawdLab такую модель можно интерпретировать следующим образом. Рабочий уровень включает текущую задачу, активную попытку выполнения, последние provider jobs, текущий план и ближайшие артефакты. Эпизодический уровень хранит историю попыток: какие варианты уже проверялись, какие замечания были получены, какие ошибки возникли и какие изменения были внесены. Долгосрочный или канонический уровень содержит принятые выводы, финальные документы, устойчивые ограничения и обобщенные уроки по задаче. В отличие от полной графовой памяти, HCC-подход можно реализовывать эволюционно: сначала явно разделить уровни памяти и правила перехода информации между ними, а затем при необходимости добавлять более сложные связи.

Отдельным направлением является представление памяти в виде карты знаний (mind map) или графа рассуждения. В работе Yilin Wen, Zifeng Wang и Jimeng Sun "MindMap" предлагается использовать граф знаний как опору для рассуждения LLM и повышения прозрачности ответа [3]. Технически такой подход не обязательно является основным хранилищем памяти. Скорее он может работать как навигационный слой: система строит карту ключевых сущностей и связей, чтобы агенту было проще ориентироваться в сложной задаче и объяснять ход рассуждения. Для ClawdLab это направление может быть полезно как следующий этап развития интерфейса памяти, например для визуализации связей между гипотезами, попытками и критикой.

С учетом этих вариантов для экспериментальной реализации выбирается промежуточный подход: структурированная артефактная память с элементами иерархического когнитивного кэширования. Она не требует полной логической графовой модели, но фиксирует явные связи между ключевыми объектами исследования и разделяет память по уровням детализации. Минимальный набор таких связей включает:

- задача связана с исходной гипотезой и целью,
- попытка выполнения связана с задачей,
- результат связан с попыткой,
- критика связана с результатом,
- новая попытка связана с причиной пересмотра,
- итоговый вывод связан с использованными результатами и ограничениями,
- принятый вывод отделен от промежуточных артефактов и черновых попыток.

Такой подход достаточно прост для прототипа и одновременно позволяет проверить основную идею: влияет ли явное структурирование памяти на согласованность, воспроизводимость и способность системы использовать результаты предыдущих попыток. В терминах HCC рабочий уровень обеспечивает выполнение текущей попытки, эпизодический уровень хранит историю попыток и критики, а канонический уровень фиксирует принятые результаты исследования.

\subsection{Технологии циклического выполнения задач}

Второе направление доработки связано с поддержкой цикличности. Под цикличностью в данном контексте понимается способность системы не завершать задачу после одного линейного прохода, а возвращаться к предыдущим шагам, если результат оказался неполным, противоречивым или недостаточно обоснованным. Технически цикличность может реализовываться на разных уровнях: внутри одного ответа модели, внутри одного агента, между несколькими агентами или на уровне всей исследовательской задачи.

Первый уровень - циклы рассуждения внутри языковой модели. В работе Yao и соавт. "Tree of Thoughts" рассуждение представлено как поиск по дереву промежуточных мыслей (thoughts) [4]. Модель не сразу генерирует финальный ответ, а создает несколько вариантов промежуточных шагов, оценивает их и продолжает наиболее перспективные ветви. Технически такой подход требует операций генерации кандидатов, оценки состояния и выбора ветви. Его достоинство состоит в том, что он позволяет исследовать несколько альтернатив вместо одного жадного вывода. Ограничение состоит в том, что цикл остается в основном внутри reasoning-процесса модели и не обязательно фиксируется как часть внешнего состояния системы.

Второй уровень - саморефлексия агента. В работе Shinn и соавт. "Reflexion" агент после выполнения задачи получает обратную связь, формирует текстовое размышление о своих ошибках и использует его в следующей попытке [5]. В отличие от Tree of Thoughts, здесь важна не только генерация альтернатив, но и сохранение опыта предыдущего выполнения. Технически Reflexion добавляет в цикл отдельный компонент памяти, где хранятся выводы о прошлых ошибках. Такой подход хорошо подходит для задач, где можно явно оценить результат и затем улучшить следующую попытку без дообучения модели.

Третий уровень - многоагентное обсуждение и эволюция гипотез. В работе "Towards an AI Co-Scientist" Готтесман и соавт. предлагают систему, где несколько специализированных агентов повторяют фазы генерации, обсуждения, ранжирования и эволюции гипотез [6]. В такой архитектуре цикл реализуется как коллективный процесс: гипотезы не просто создаются один раз, а проходят через критику, сравнение и улучшение. Технически это ближе к научному процессу, поскольку важным становится не только финальный ответ, но и история изменений гипотезы.

Четвертый уровень - стадийная обратная связь в исследовательском процессе. В Agent Laboratory исследование разделяется на крупные этапы: обзор литературы, экспериментирование и подготовка отчета [7]. Пользователь может давать обратную связь на границах стадий, после чего система корректирует дальнейшие действия. Такой подход можно назвать стадийным циклом (stage-gated feedback): система сохраняет управляемую последовательность этапов, но допускает уточнение внутри или между ними. Для ClawdLab этот вариант особенно близок, поскольку платформа также предполагает участие человека и наличие управляющей роли.

Пятый уровень - долгосрочный исследовательский цикл. В DeepScientist процесс описан как повторение фаз "гипотеза - проверка - анализ" (hypothesize - verify - analyze), где последующие шаги опираются на накопленную память находок [8]. Технически такая модель требует не только запуска повторных попыток, но и сохранения результатов предыдущих проверок в форме, пригодной для выбора следующего направления. Иными словами, цикличность здесь тесно связана с памятью: без накопления результатов система не сможет понять, какие гипотезы уже проверялись и какие направления стоит развивать дальше.

Для ClawdLab прямой перенос сложных схем, таких как полноценный турнир гипотез, многоагентная эволюция или динамическая перестройка архитектуры, не является необходимым для первого прототипа. Более подходящей является модель управляемых попыток выполнения задачи (TaskAttempt или TaskRun). В этой модели верхнеуровневая задача сохраняет роль исследовательской цели, а каждая попытка фиксирует один конкретный проход решения: план, входные данные, действия агентов, полученные результаты, критику и решение по итогам проверки.

Такой подход позволяет реализовать цикл без разрушения существующего жизненного цикла задач. Если попытка завершилась недостаточно качественным результатом, система создает новую попытку, явно связывая ее с причиной пересмотра. Например, причиной может быть недостаточная полнота источников, противоречие в аргументации, слабая обоснованность вывода или замечание критика. При этом предыдущая попытка не удаляется и не перезаписывается: она остается частью истории исследования и может использоваться для анализа ошибок.

С технической точки зрения сущность попытки должна хранить как минимум следующие данные:

- номер попытки и ссылку на исходную задачу,
- краткий план выполнения,
- использованные входные данные и источники,
- связанные provider jobs и артефакты,
- результат попытки,
- замечания критика или управляющего агента,
- решение по итогам проверки,
- причину повторного запуска, если создается новая попытка,
- отличия новой попытки от предыдущей.

Минимальный цикл для эксперимента включает следующие шаги:

- постановка задачи,
- выполнение первой попытки,
- фиксация результата и артефактов,
- критика результата,
- принятие решения о завершении или повторе,
- выполнение следующей попытки с учетом критики,
- сравнение результата новой попытки с предыдущей.

Важным элементом такой модели являются правила остановки. Цикл не должен быть бесконечным: необходимо заранее задать максимальное число попыток, условия принятия результата и случаи, когда задача завершается как недостаточно подтвержденная. Это позволяет сохранить управляемость и сопоставимость экспериментов.

Таким образом, для экспериментальной реализации выбирается не полный перенос сложных reasoning-алгоритмов, а системный цикл на уровне исследовательской задачи. Такая модель проверяет именно ту часть цикличности, которая важна для ClawdLab: способность не просто создать новый диалог или новую задачу, а сохранить историю попыток внутри одного исследовательского процесса и использовать эту историю для улучшения следующего шага.

\subsection{Инструменты оркестрации и реализации}

В качестве базовой платформы рассматривается ClawdLab, построенный поверх OpenClaw. OpenClaw предоставляет инфраструктурные возможности для подключения агентов, каналов, сессий, инструментов и плагинов [9]. ClawdLab добавляет к этой основе лабораторную организацию: роли, управление, критику, артефакты и проверку результатов [10].

Для экспериментального прототипа важно не заменить эту архитектуру, а расширить ее минимальным набором новых сущностей и связей. На концептуальном уровне необходимы следующие элементы:

- сущность попытки выполнения задачи,
- связи между попыткой, артефактами, критикой и итоговым решением,
- слой памяти, позволяющий извлекать релевантные артефакты предыдущих попыток,
- механизм принятия решения о повторе или завершении цикла.

В качестве инструментального слоя могут использоваться существующие возможности ClawdLab: внешние агенты, provider jobs, документы, обсуждения и журналы активности. Это позволяет сохранить сопоставимость между базовой и улучшенной конфигурациями. Различие между ними должно заключаться не в разных моделях или разных внешних инструментах, а именно в наличии структурированной памяти и циклической логики выполнения.

\subsection{Выбор подходов для экспериментальной реализации}

По результатам обзора для дальнейшей экспериментальной реализации выбираются следующие решения.

Для памяти выбирается структурированная артефактная модель (HCC). Она является более простой, чем полноценная графовая когнитивная память, но сохраняет ключевую идею: промежуточные результаты должны быть связаны с гипотезами, попытками, критикой и итоговыми выводами.

Для цикличности выбирается модель управляемых попыток выполнения задачи. Она позволяет реализовать возврат к предыдущим этапам без полной перестройки архитектуры ClawdLab и без изменения принципа подключения внешних агентов.

Динамический выбор моделей не включается в экспериментальную реализацию. Он остается перспективным направлением, но в текущей постановке не соответствует границам системы, поскольку ClawdLab не управляет внутренним inference-контуром внешних агентов.


Источники:

1. Graph-Native Cognitive Memory for AI Agents - [https://arxiv.org/abs/2603.17244](https://arxiv.org/abs/2603.17244)
2. Toward Ultra-Long-Horizon Agentic Science: Cognitive Accumulation for Machine Learning Engineering - [https://arxiv.org/abs/2601.10402](https://arxiv.org/abs/2601.10402)
3. MindMap: Knowledge Graph Prompting Sparks Graph of Thoughts in Multi-Hop Reasoning - [https://arxiv.org/abs/2308.09729](https://arxiv.org/abs/2308.09729)
4. Tree of Thoughts: Deliberate Problem Solving with Large Language Models - [https://arxiv.org/abs/2305.10601](https://arxiv.org/abs/2305.10601)
5. Reflexion: Language Agents with Verbal Reinforcement Learning - [https://arxiv.org/abs/2303.11366](https://arxiv.org/abs/2303.11366)
6. Towards an AI Co-Scientist - [https://arxiv.org/abs/2502.18864](https://arxiv.org/abs/2502.18864)
7. Agent Laboratory: Using LLM Agents as Research Assistants - [https://arxiv.org/abs/2501.04227](https://arxiv.org/abs/2501.04227)
8. DeepScientist: Advancing Frontier-Pushing Scientific Findings Progressively - [https://arxiv.org/abs/2509.26603](https://arxiv.org/abs/2509.26603)
9. OpenClaw documentation - [https://docs.openclaw.ai/](https://docs.openclaw.ai/)
10. From Agent-Only Social Networks to Autonomous Scientific Research: Lessons from OpenClaw and Moltbook, and the Architecture of ClawdLab and Beach.Science - [https://arxiv.org/abs/2602.19810](https://arxiv.org/abs/2602.19810)


\section{Экспериментальные данные и задачи}

\subsection{Выбор типов задач для оценки}

Экспериментальная проверка проводится не на всех типах задач, предусмотренных в ClawdLab, а на двух типах, наиболее связанных с целью данной работы: задачах анализа (analysis) и задачах глубокого исследования (deep research) [1]. Такой выбор позволяет не выходить за рамки классификации задач, предложенной в статье о ClawdLab, и одновременно проверить обе выбранные доработки: структурированную память и циклическое выполнение.

В рамках данной работы лаборатория используется как рабочее пространство для исследовательского вопроса, однако экспериментальная проверка ограничивается отдельными задачами типов analysis и deep research. Остальные типы задач ClawdLab, такие как literature review, critique и synthesis, могут использоваться как вспомогательные стадии внутри выполнения, но не являются самостоятельными объектами сравнения.

Задачи анализа подходят для проверки того, насколько система умеет работать с заданным набором источников, выделять аргументы, сопоставлять подходы и формировать обоснованный вывод. Задачи глубокого исследования требуют более длинного горизонта работы: система должна не только проанализировать входные материалы, но и уточнить гипотезу, выявить ограничения, сформировать промежуточные выводы и при необходимости вернуться к предыдущим шагам. Поэтому первый тип задач больше проверяет качество структурирования и согласованность вывода, а второй - устойчивость процесса при более сложной исследовательской логике.

В рамках эксперимента единицей сравнения является не абстрактная "производительность системы в целом", а результат выполнения одинаковой научной задачи в разных конфигурациях ClawdLab. Иными словами, каждая конфигурация получает один и тот же тип задачи, одну и ту же гипотезу, одинаковые входные материалы и одинаковые ограничения. После этого сравниваются итоговые артефакты и процесс их получения: качество вывода, полнота учета источников, согласованность с промежуточными результатами, использование критики, количество повторных попыток и воспроизводимость хода работы.

Такой подход позволяет оценивать систему не только по финальному тексту, но и по тому, насколько управляемо она пришла к результату. Это важно для ClawdLab, поскольку предлагаемая доработка касается не замены языковой модели и не ускорения отдельного вызова, а организации исследовательского процесса. Поэтому в эксперименте сравниваются результаты выполнения задач, а показатели системы рассматриваются как дополнительные метрики процесса.

\subsection{Задачи анализа научной информации}

Для задач анализа научной информации подготавливается несколько исследовательских гипотез, каждая из которых может быть проверена на фиксированном наборе материалов. Входными данными для такой задачи являются формулировка гипотезы, набор статей или технических описаний, а также ожидаемый формат результата. Система должна выделить релевантные аргументы, сопоставить позиции источников, указать подтверждения и ограничения, а затем сформировать структурированный аналитический вывод.

Примером такой задачи может быть проверка гипотезы о том, что структурированная память снижает потерю контекста в длинных агентных процессах. Для этой гипотезы система получает набор источников о памяти агентов, графовых представлениях, иерархическом кэшировании и научных агентных системах. Результатом должен стать аналитический отчет, в котором указано, какие источники поддерживают гипотезу, какие показывают ограничения подхода и какие выводы можно использовать для проектирования ClawdLab.

Для каждой задачи анализа фиксируются:

- формулировка проверяемой гипотезы,
- набор входных источников,
- ожидаемый тип результата,
- критерии полноты и качества вывода,
- ограничения на число попыток выполнения.

Использование нескольких гипотез позволяет уменьшить зависимость результата от одной удачной или неудачной постановки. При этом все конфигурации системы выполняют один и тот же набор задач, что обеспечивает сопоставимость результатов.

\subsection{Задачи глубокого исследования}

Задачи глубокого исследования отличаются от задач анализа тем, что требуют более развернутой исследовательской траектории. В них система должна не только обработать известный набор материалов, но и сформировать более целостное понимание проблемы: уточнить исходную гипотезу, определить недостающие сведения, предложить возможные направления проверки, провести критику промежуточного результата и синтезировать итоговый вывод.

В таких задачах входная постановка может быть менее закрытой, чем в задачах анализа. Например, вместо требования проверить одну конкретную гипотезу система может получить исследовательский вопрос: какие архитектурные механизмы наиболее важны для повышения воспроизводимости мультиагентной научной системы. В этом случае результат должен включать не только ответ, но и ход рассуждения: какие подзадачи были выделены, какие источники использованы, какие промежуточные выводы были пересмотрены после критики и почему итоговое решение считается обоснованным.

Для задач глубокого исследования особенно важны обе доработки. Структурированная память должна помогать системе сохранять связь между исходным вопросом, промежуточными выводами, замечаниями критика и итоговым синтезом. Циклическая модель должна позволять не завершать задачу после первого результата, если критика показывает недостаточную полноту, противоречие или слабую аргументацию. Поэтому данный тип задач используется как более сложный сценарий проверки целевой архитектуры.

Для эксперимента планируется подготовить несколько задач глубокого исследования, каждая из которых задается через:

- исследовательский вопрос или исходную гипотезу,
- начальный набор источников,
- допустимый формат дополнительных уточнений,
- ожидаемый итоговый артефакт,
- критерии оценки глубины, согласованности и учета критики.

\subsection{Подготовка входных данных и эталонных результатов}

Входные данные для эксперимента должны быть одинаковыми для всех сравниваемых конфигураций. Для каждой задачи формируется карточка эксперимента, включающая тип задачи, формулировку гипотезы или исследовательского вопроса, список исходных материалов, ограничения на выполнение и критерии оценки. Такой формат позволяет запускать одну и ту же задачу в базовой и улучшенных конфигурациях без изменения содержания постановки.

Для задач анализа основным входом является фиксированный набор источников. Это позволяет оценить, насколько система корректно использует уже предоставленную информацию. Для задач глубокого исследования входной набор также фиксируется, но допускается более сложная работа с промежуточными выводами и критикой. При этом эксперимент не должен превращаться в проверку качества внешнего веб-поиска, поэтому дополнительные источники, если они разрешены, должны фиксироваться и учитываться одинаково для всех конфигураций.

Эталонный результат в данном эксперименте не обязательно является единственным правильным ответом. Для исследовательских задач часто невозможно заранее задать полностью формальный ответ, особенно если речь идет о сравнении подходов или уточнении гипотез. Поэтому эталон используется как набор критериев и ожидаемых элементов результата: какие источники должны быть учтены, какие аргументы являются ключевыми, какие ограничения необходимо указать и какие противоречия нельзя игнорировать.

Таким образом, подготовка данных включает три уровня:

- набор задач анализа и глубокого исследования,
- карточки задач с входными материалами и ограничениями,
- экспертные критерии оценки результатов и процесса выполнения.

Источники:

1. From Agent-Only Social Networks to Autonomous Scientific Research: Lessons from OpenClaw and Moltbook, and the Architecture of ClawdLab and Beach.Science - [https://arxiv.org/abs/2602.19810](https://arxiv.org/abs/2602.19810)

\section{Экспериментальные конфигурации}

\subsection{Базовая конфигурация ClawdLab}

Базовая конфигурация соответствует текущей логике ClawdLab без добавления структурированной памяти исследования и без формализованного цикла попыток. В ней задача выполняется как управляемый агентный процесс: система принимает постановку, распределяет работу между агентами, сохраняет обсуждения и артефакты, а затем формирует результат. При этом связи между гипотезой, промежуточными выводами, критикой и итоговым ответом остаются в основном неявными.

Эта конфигурация используется как контрольная. Она показывает, какого качества можно достичь при существующей организации процесса и какие ограничения проявляются при выполнении задач анализа и глубокого исследования. Все остальные конфигурации сравниваются с ней на одном и том же наборе задач.

\subsection{Конфигурация со структурированной памятью}

Во второй конфигурации к базовой системе добавляется структурированная память исследования. Она фиксирует ключевые объекты задачи и связи между ними: исходную гипотезу, попытку выполнения, использованные источники, промежуточные выводы, замечания критика, итоговый результат и ограничения. Основная цель этой конфигурации - проверить, улучшает ли явное связывание артефактов согласованность результата и способность системы использовать уже полученный контекст.

В этой конфигурации процесс выполнения остается преимущественно линейным: задача не получает отдельного механизма повторных попыток. Это позволяет изолированно оценить вклад памяти. Если результат улучшается по полноте, согласованности или воспроизводимости, но не улучшается по учету критики, это может означать, что одной памяти недостаточно без циклической логики пересмотра.

\subsection{Конфигурация с циклическим процессом}

Третья конфигурация добавляет к базовой системе механизм управляемых попыток выполнения задачи. После получения первичного результата система проводит критику и принимает решение: завершить задачу или создать новую попытку с учетом замечаний. Каждая попытка имеет собственный результат, причину запуска и отличие от предыдущего прохода.

В этой конфигурации структурированная память в полном виде не добавляется. Сохраняется только минимальная информация, необходимая для работы цикла: номер попытки, результат, критика и причина повторного запуска. Это позволяет проверить вклад цикличности отдельно от более развитого слоя памяти. Ожидается, что такая конфигурация может улучшить качество проработки спорных или неполных выводов, но без структурированной памяти часть связей между артефактами все еще может оставаться неявной.

\subsection{Итоговая конфигурация с обеими доработками}

Итоговая конфигурация объединяет структурированную память и циклический процесс. В ней каждая попытка выполнения задачи сохраняется как часть истории исследования, а ее результаты связываются с гипотезой, источниками, критикой и итоговым решением. Если создается новая попытка, система фиксирует, какая проблема исправляется и какие элементы предыдущего результата должны быть учтены.

Эта конфигурация соответствует целевому образу решения TO BE, сформулированному в первой главе. Она должна показать, дает ли совместное использование памяти и цикличности больший эффект, чем каждая доработка по отдельности. Особенно важна проверка на задачах глубокого исследования, где потеря контекста и необходимость пересмотра промежуточных выводов проявляются сильнее, чем в более простых аналитических задачах.

\subsection{Ablation-сравнение}

Для оценки вклада отдельных модулей используется ablation-сравнение, то есть последовательное сравнение конфигураций с разным набором доработок. В эксперименте рассматриваются четыре варианта:

- базовая конфигурация,
- конфигурация только со структурированной памятью,
- конфигурация только с циклическим процессом,
- итоговая конфигурация со структурированной памятью и цикличностью.

Такое сравнение позволяет ответить на три вопроса. Во-первых, дает ли структурированная память самостоятельный эффект по сравнению с базовой конфигурацией. Во-вторых, дает ли циклический процесс самостоятельный эффект. В-третьих, усиливают ли эти механизмы друг друга при совместном использовании.

Для корректного сравнения во всех конфигурациях должны быть зафиксированы одинаковые внешние условия: набор задач, входные материалы, роли агентов, используемые модели, доступные инструменты, лимиты выполнения и критерии оценки. Это особенно важно, поскольку в данной работе не проверяется динамический выбор моделей. Следовательно, различия в результатах должны объясняться не заменой модели или инструмента, а именно изменением архитектуры исследовательского процесса.

\section{Методика оценки}

\subsection{Общая логика валидации}

Цель оценки состоит не в том, чтобы показать, что система "в целом стала лучше", а в том, чтобы проверить конкретное архитектурное предположение: структурированная память и управляемая цикличность могут улучшить качество, согласованность и воспроизводимость выполнения научных задач в ClawdLab. Поэтому методика строится как сравнительный эксперимент с фиксированными условиями: одни и те же задачи, источники, роли агентов, модели, лимиты и критерии оценки запускаются в нескольких конфигурациях системы.

В терминах дизайна экспериментов по Фишеру заранее фиксируются шесть элементов: исследовательский вопрос, данные, метрики, оцениваемое свойство системы, дизайн запуска и критерий принятия решения [8]. Это важно, потому что для LLM- и агентных систем легко подобрать удобную метрику уже после получения результата. Чтобы снизить риск такой подгонки, ожидаемые эффекты и критерии приемки должны быть определены до запуска эксперимента.

Нулевая гипотеза $H_0$ формулируется следующим образом: добавление структурированной памяти и циклического выполнения не приводит к значимому улучшению качества, согласованности и воспроизводимости по сравнению с базовой конфигурацией ClawdLab. Альтернативная гипотеза $H_A$ состоит в том, что хотя бы одна из доработок дает измеримый положительный вклад, а итоговая конфигурация с обеими доработками показывает лучший результат по основным метрикам без недопустимого ухудшения стоимости и времени выполнения.

Для проверки используется не одна метрика, а несколько групп показателей. Такой подход соответствует идее многоосевой оценки, применяемой в HELM: качество языковой системы нельзя описывать только точностью, поскольку важны также устойчивость, эффективность, безопасность и trade-off между ними [3]. Для данной работы эта идея адаптируется к агентной научной системе: оценивается не только финальный текст, но и процесс его получения, связь с артефактами, способность учитывать критику и возможность восстановить ход выполнения.

\subsection{Оценка итогового результата}

Первая группа метрик оценивает качество итогового артефакта, полученного после выполнения задачи analysis или deep research. Поскольку для исследовательских задач часто нет единственного правильного ответа, вместо бинарной проверки используется экспертная рубрика. Такой подход близок к оценке Agent Laboratory, где исследователи оценивали полезность, качество отчета, качество эксперимента и пригодность результата для продолжения работы [1].

В рамках данной работы итоговый результат оценивается по следующим критериям:

- полнота ответа: учтены ли основные аспекты исследовательского вопроса,
- обоснованность: связаны ли выводы с источниками, данными или артефактами,
- фактическая корректность: отсутствуют ли неподтвержденные утверждения и искажения источников,
- аналитическая глубина: есть ли сравнение подходов, ограничений и альтернативных интерпретаций,
- согласованность: не противоречат ли итоговые выводы промежуточным результатам и критике,
- полезность для продолжения исследования: можно ли на основе результата принять дальнейшее решение.

Каждый критерий может оцениваться по дискретной шкале, например от 1 до 5. Для уменьшения субъективности к шкале добавляются описания уровней и примеры оценок. Итоговая оценка качества не заменяет анализ отдельных критериев: если система набирает высокий общий балл, но плохо учитывает критику или источники, это должно быть видно в разборе.

Возможна дополнительная автоматизированная оценка с помощью LLM-as-a-Judge, но она не должна быть единственным источником вывода. Современные работы по LLM-судьям показывают, что такие оценки требуют калибровки на человеческой разметке и могут быть смещены из-за неоднозначности шкалы, предпочтения более длинных ответов или различий между интерпретациями рецензентов [6, 7]. Поэтому LLM-судья используется как вспомогательный инструмент: промпт судьи фиксируется заранее, температура устанавливается близкой к нулю, примеры перемешиваются, а часть оценок проверяется человеком.

Практические руководства по оценке LLM-систем также рекомендуют разделять автоматические бенчмарки, человеческую оценку и LLM-as-a-Judge, поскольку эти подходы имеют разные ограничения по стоимости, скорости, воспроизводимости и качеству [4]. Для части критериев могут использоваться готовые LLM-evaluation подходы, например проверка полноты ответа, фактической заземленности, релевантности контекста и устойчивости к prompt injection [5], но в данной работе они должны быть адаптированы к исследовательским артефактам ClawdLab.

\subsection{Оценка артефактов и исследовательского процесса}

Вторая группа метрик связана с тем, что ClawdLab оценивается не только как генератор текста, но и как среда организации исследовательской работы. Поэтому важным объектом оценки являются артефакты: карточка задачи, исходная гипотеза, использованные источники, промежуточные выводы, результаты попыток, замечания критика, решения о повторном запуске и итоговый документ.

Для базовой конфигурации артефакты могут существовать как отдельные сообщения, файлы или записи журнала. В целевой конфигурации они должны быть связаны между собой. Поэтому оцениваются следующие показатели:

- полнота трассы выполнения: можно ли восстановить, какие шаги были выполнены и в каком порядке,
- покрытие артефактов: созданы ли ожидаемые артефакты для данного типа задачи,
- связность артефактов: указаны ли связи между гипотезой, попыткой, источниками, критикой и выводом,
- происхождение вывода: понятно ли, из каких источников и промежуточных результатов получен итоговый ответ,
- учет критики: отражены ли замечания критика в следующей попытке или в итоговом результате,
- фиксация ограничений: сохранены ли известные слабые места, ошибки и нерешенные вопросы.

Именно эта часть оценки проверяет вклад структурированной памяти. Если итоговый текст формально выглядит качественным, но невозможно понять, откуда взялись выводы и какие замечания были учтены, то ожидаемый эффект от памяти считается подтвержденным только частично. Напротив, улучшенная конфигурация должна показывать более явную и воспроизводимую связь между элементами исследования.

\subsection{Метрики цикличности и использования критики}

Третья группа метрик проверяет вклад управляемого цикла выполнения. Для этого оценивается не только наличие повторных попыток, но и их осмысленность. Повторный запуск сам по себе не является улучшением, если система просто генерирует новый вариант ответа без понятной причины.

Для каждой задачи фиксируются:

- количество попыток выполнения,
- причина создания новой попытки,
- тип исправляемой проблемы: неполнота, противоречие, ошибка источника, слабая аргументация, недостаточная проверка,
- доля замечаний критика, которые были учтены в следующей попытке,
- изменение качества результата между попытками,
- наличие регресса, когда новая попытка ухудшает ранее корректный результат.

В задачах научных агентов такая оценка особенно важна, потому что современные бенчмарки показывают падение качества при длинном горизонте действий. Например, SciAgentGym оценивает не только финальный успех, но и способность агента выполнять многошаговые научные workflows, авторы показывают, что при росте длины траектории качество инструментального выполнения заметно падает [2]. Для ClawdLab это означает, что цикличность должна оцениваться как управляемый механизм исправления ошибок, а не как бесконтрольное увеличение числа шагов.

\subsection{Метрики воспроизводимости и устойчивости}

Воспроизводимость в данной работе понимается как возможность восстановить ход выполнения задачи и получить сопоставимый результат при повторном запуске в тех же условиях. Для агентных систем это особенно сложно, поскольку результат зависит от модели, промпта, порядка сообщений, доступных инструментов и состояния памяти.

Оцениваются следующие показатели:

- полнота протокола запуска: сохранены модель, версия промптов, входные данные, ограничения, параметры и время запуска,
- восстановимость траектории: можно ли по логам и артефактам повторить основные шаги выполнения,
- стабильность результата: насколько сильно меняется оценка качества при повторных запусках,
- стабильность связей памяти: сохраняются ли связи между гипотезами, попытками, критикой и выводами,
- отсутствие утечек: не попадает ли в входные данные информация из эталонной оценки или результатов других конфигураций.

Для проверки устойчивости набор задач должен включать не только "хорошие" примеры, но и сложные случаи. Поэтому датасет целесообразно разделить на три части: golden set с типовыми экспертными примерами, adversarial set со сложными или провокационными постановками и regression set с примерами, где система ранее ошибалась. Такая структура позволяет показать не только среднее качество, но и поведение системы на слабых местах.

Отдельно проверяется, что улучшенная конфигурация не деградирует на нецелевых или пограничных задачах. Например, если задача выходит за рамки analysis или deep research, система должна корректно обозначить ограничение, а не уверенно выполнять неподходящий сценарий. Для агентной системы также важно отслеживать устойчивость к шуму во входных данных и попыткам prompt injection, особенно если агенты работают с внешними источниками и инструментами.

\subsection{Метрики вычислительной эффективности}

Поскольку предлагаемая доработка добавляет память и дополнительные попытки выполнения, она может повышать качество ценой роста затрат. Поэтому вместе с качеством фиксируются инженерные метрики:

- общее время выполнения задачи,
- число вызовов агентов и внешних инструментов,
- число токенов или примерная стоимость выполнения,
- количество созданных артефактов,
- число попыток выполнения,
- доля неудачных или прерванных запусков.

Эти метрики не являются основными целевыми метриками, но они нужны для честного анализа trade-off и для принятия решения по $H_A$. Например, если итоговая конфигурация улучшает качество на небольшую величину, но увеличивает стоимость в несколько раз, такой результат трудно считать практически полезным. Напротив, умеренное увеличение времени может быть приемлемым, если система становится заметно более воспроизводимой и лучше учитывает критику.

\subsection{Критерии принятия гипотезы}

Решение по гипотезам принимается относительно формулировок $H_0$ и $H_A$, заданных выше. Нулевая гипотеза $H_0$ не отклоняется, если улучшенные конфигурации не показывают устойчивого преимущества перед базовой конфигурацией или если наблюдаемое улучшение объясняется единичными примерами, ростом стоимости выполнения или изменением внешних условий эксперимента.

Нулевая гипотеза $H_0$ отклоняется в пользу альтернативной гипотезы $H_A$, если итоговая конфигурация со структурированной памятью и цикличностью показывает улучшение по заранее выбранным целевым метрикам по сравнению с базовой конфигурацией, а ablation-сравнение показывает измеримый вклад хотя бы одной из доработок. Основными целевыми метриками являются:

- качество итогового результата,
- согласованность результата с источниками, критикой и промежуточными выводами,
- полнота и связность артефактов,
- воспроизводимость траектории выполнения,
- учет критики в последующих попытках.

Для количественных или порядковых оценок используется парное сравнение: одна и та же задача сравнивается в разных конфигурациях. Если размер выборки достаточен, для оценки устойчивости результата применяются доверительные интервалы, бутстрап или непараметрические тесты для связанных выборок. Если выборка мала, вывод не должен опираться только на p-value, в этом случае основную роль играет качественная абляция, показывающая, какой именно компонент системы дал вклад.

Анализ проводится в несколько шагов. Сначала сравнивается базовая конфигурация и конфигурация только со структурированной памятью. Затем сравнивается базовая конфигурация и конфигурация только с цикличностью. После этого оценивается итоговая конфигурация с обеими доработками. Такой ablation-дизайн позволяет показать не только факт улучшения, но и источник эффекта: память, цикл или их совместное действие.

Для отклонения $H_0$ недостаточно, чтобы итоговый отчет выглядел лучше в одном удачном примере. Нужно показать, что улучшение повторяется на нескольких задачах, не сопровождается сильной деградацией стоимости и объясняется именно архитектурными изменениями. Если улучшение наблюдается только по качеству текста, но не по связности артефактов и воспроизводимости, результат интерпретируется как частичное свидетельство в пользу $H_A$. Если структурированная память улучшает трассируемость, но не влияет на итоговое качество, это также фиксируется как частичный результат и основание для дальнейшей доработки.

\subsection{Протокол проведения экспериментов}

Эксперимент проводится по единому протоколу. Сначала формируется набор задач analysis и deep research, включающий golden, adversarial и regression-примеры. Для каждой задачи заранее создается карточка: формулировка, входные материалы, ожидаемый тип результата, ограничения, критерии оценки и допустимое число попыток.

Затем каждая задача запускается во всех экспериментальных конфигурациях: базовой, с памятью, с циклом и с обеими доработками. Порядок запусков перемешивается, чтобы снизить влияние порядка оценки. Для всех запусков фиксируются версии моделей, промптов, инструментов, параметры выполнения и лимиты.

После выполнения сохраняются все итоговые и промежуточные артефакты. Оценщик получает не только финальный результат, но и структурированный пакет выполнения: карточку задачи, список источников, попытки, критику, решения о пересмотре и итоговый документ. При этом для экспертной оценки желательно скрывать название конфигурации, чтобы рецензент не знал, какой вариант системы он оценивает.

Далее проводится оценка по рубрикам. Часть примеров размечается человеком и используется для калибровки LLM-судьи. После этого LLM-судья может применяться для масштабирования оценки, но спорные и худшие случаи дополнительно проверяются вручную. Для каждого типа ошибок формируется разбор: где система потеряла контекст, где неправильно использовала источник, где не учла критику, где цикл ухудшил результат или привел к лишним действиям.

Итоговый отчет по эксперименту должен содержать не только средние значения метрик, но и анализ ошибок. Для защиты это принципиально важно: демонстрация слабых мест системы показывает, что эксперимент не сводится к подбору удачных примеров, а действительно проверяет границы применимости решения.

Источники:

1. Agent Laboratory: Using LLM Agents as Research Assistants - [https://arxiv.org/abs/2501.04227](https://arxiv.org/abs/2501.04227)
2. SciAgentGym: Benchmarking Multi-Step Scientific Tool-use in LLM Agents - [https://arxiv.org/abs/2602.12984](https://arxiv.org/abs/2602.12984)
3. Holistic Evaluation of Language Models - [https://arxiv.org/abs/2211.09110](https://arxiv.org/abs/2211.09110)
4. Hugging Face LLM Evaluation Guidebook - [https://github.com/huggingface/evaluation-guidebook](https://github.com/huggingface/evaluation-guidebook)
5. UpTrain predefined evaluations documentation - [https://docs.uptrain.ai/predefined-evaluations/overview](https://docs.uptrain.ai/predefined-evaluations/overview)
6. Validating LLM-as-a-Judge Systems under Rating Indeterminacy - [https://www.microsoft.com/en-us/research/publication/validating-llm-as-a-judge-systems-under-rating-indeterminacy/](https://www.microsoft.com/en-us/research/publication/validating-llm-as-a-judge-systems-under-rating-indeterminacy/)
7. Human-anchored longitudinal comparison of generative AI with a bias-calibrated LLM-as-judge - [https://journals.plos.org/plosone/article?id=10.1371/journal.pone.0339920](https://journals.plos.org/plosone/article?id=10.1371/journal.pone.0339920)
8. Fisher R. A. The Design of Experiments. Oliver and Boyd, 1935.

\section{Ожидаемый Proof-of-Concept}

\subsection{Целевая архитектура экспериментального прототипа}

Ожидаемый Proof-of-Concept в данной работе представляет собой не отдельную новую платформу, а экспериментальное расширение ClawdLab, предназначенное для проверки двух архитектурных доработок: структурированной памяти исследования и управляемой цикличности выполнения задач. Прототип должен сохранить текущую роль ClawdLab как слоя координации внешних агентов, но добавить к нему более явное представление исследовательского состояния.

На уровне целевой архитектуры прототип включает несколько основных компонентов.

Первый компонент - существующее рабочее пространство лаборатории. Лаборатория задает общий исследовательский контекст: научный вопрос, роли агентов, обсуждения, документы, артефакты и журналы активности. В рамках эксперимента лаборатория рассматривается как контейнер для одного исследовательского направления, внутри которого запускаются отдельные задачи типов analysis и deep research.

Второй компонент - расширенная модель задачи. Сущность Task сохраняет роль исследовательской цели: она описывает, что именно необходимо проанализировать, какую гипотезу проверить или какой исследовательский вопрос проработать. При этом задача не должна смешиваться с конкретным проходом выполнения. Для этого вводится отдельный уровень попыток выполнения.

Третий компонент - сущность TaskAttempt или TaskRun. Она фиксирует один конкретный проход решения задачи: номер попытки, план выполнения, использованные источники, связанные provider jobs, созданные артефакты, результат, замечания критика, решение по итогам проверки и отличие от предыдущей попытки. Такой подход позволяет не менять базовый жизненный цикл задачи, но при этом сделать цикличность явной и машинно-обрабатываемой.

Четвертый компонент - структурированная память исследования. Для прототипа выбирается не полноценная графовая память с формальным пересмотром убеждений, а более простая многоуровневая модель, близкая к Hierarchical Cognitive Caching [1]. В ней можно выделить три уровня:

- рабочий уровень: текущая задача, активная попытка, последние действия агентов и ближайшие артефакты,
- эпизодический уровень: история попыток, критика, ошибки, промежуточные выводы и причины повторных запусков,
- канонический уровень: принятые выводы, финальные документы, ограничения и устойчивые результаты исследования.

Пятый компонент - слой связей между артефактами. В базовой конфигурации ClawdLab уже хранит обсуждения, документы, provider jobs и журналы активности, однако эти элементы не всегда образуют явную исследовательскую структуру. В прототипе должны фиксироваться связи вида: задача связана с гипотезой, попытка связана с задачей, результат связан с попыткой, критика связана с результатом, новая попытка связана с причиной пересмотра, а итоговый вывод связан с использованными источниками и ограничениями.

Таким образом, целевая архитектура PoC может быть представлена как расширение существующих сущностей ClawdLab, а не как их замена. Главное изменение состоит в том, что исследовательский процесс становится представлен не только через сообщения и документы, но и через структуру: Task, TaskAttempt, артефакты, критику, решения и уровни памяти.

\subsection{Ожидаемые сценарии работы}

Прототип должен поддерживать два основных экспериментальных сценария, соответствующих выбранным типам задач ClawdLab.

Первый сценарий - задача анализа. Пользователь или главный исследователь создает лабораторию для определенного научного вопроса и внутри нее формулирует задачу типа analysis. Для задачи задаются исходная гипотеза, набор источников, ожидаемый формат результата и ограничения выполнения. Система запускает первую попытку, в рамках которой агент анализирует материалы, выделяет аргументы, формирует промежуточные выводы и создает аналитический артефакт.

После выполнения попытки результат передается на критику. Критик или главный исследователь фиксирует замечания: например, неполный учет источников, противоречие в аргументации, слабую связь между выводом и доказательствами или отсутствие ограничений. Если результат признается достаточным, он переходит в канонический уровень памяти как принятый вывод. Если результат требует доработки, создается новая попытка, связанная с предыдущей и с причиной пересмотра.

Второй сценарий - задача глубокого исследования. В этом случае входная постановка может быть шире: не только проверить одну гипотезу, но и проработать исследовательский вопрос, выделить возможные направления проверки, сопоставить подходы и подготовить более развернутый вывод. Такой сценарий требует более длинной траектории: несколько промежуточных выводов, критику, уточнение исходной постановки и итоговый синтез.

Для deep research особенно важна связка памяти и цикличности. Структурированная память должна сохранять, какие подзадачи были выделены, какие источники использованы, какие промежуточные выводы были получены и какие замечания появились после критики. Циклическая модель должна позволять вернуться к предыдущей попытке неформально не "с нуля", а как к документированному этапу исследования: с указанием причины, измененного плана и ожидаемого улучшения.

Ожидаемый результат работы прототипа включает не только финальный текстовый документ, но и пакет исследовательских артефактов:

- карточку задачи,
- список исходных источников,
- одну или несколько попыток выполнения,
- результаты provider jobs,
- промежуточные выводы,
- замечания критика,
- решения о принятии результата или повторном запуске,
- итоговый отчет,
- связи между задачей, попытками, критикой, источниками и выводами.

Именно этот пакет артефактов используется в дальнейшем для оценки качества, воспроизводимости и вклада архитектурных доработок.

\subsection{Ограничения прототипа}

Ожидаемый прототип имеет несколько принципиальных ограничений.

Во-первых, он не должен реализовывать полный автономный цикл научного открытия. Его задача состоит в проверке двух конкретных механизмов внутри ClawdLab: структурированной памяти и циклического выполнения. Поэтому прототип ограничивается задачами analysis и deep research и не претендует на автоматизацию всех типов научной деятельности.

Во-вторых, в прототип не включается динамический выбор языковых моделей. ClawdLab работает с внешними агентами, которые имеют собственные модели и исполнительные контуры. Платформа может передавать агентам контекст, роль, задачу и артефакты, но не управляет напрямую тем, какую модель внешний агент использует внутри своего runtime. Поэтому выбор моделей фиксируется как условие эксперимента, а не как часть разрабатываемой архитектуры.

В-третьих, в рамках PoC не реализуется полноценная graph-native memory с формальной семантикой поддержки, опровержения и пересмотра убеждений [2]. Такой подход является перспективным, но для первого прототипа он избыточен: он потребовал бы отдельной онтологии, правил версионирования и более сложного графового хранилища. В данной работе используется более легкая артефактно-ориентированная память с элементами иерархического кэширования.

В-четвертых, mind-map или графовая визуализация знания рассматривается как возможное развитие, но не как обязательная часть экспериментального прототипа [3]. В PoC достаточно зафиксировать машинно-обрабатываемые связи между задачами, попытками, источниками, критикой и выводами. Визуальное представление этих связей может быть добавлено позже как интерфейсный слой.

В-пятых, цикличность ограничивается заранее заданными правилами остановки: максимальным числом попыток, решением главного исследователя, отсутствием улучшения или превышением вычислительного бюджета. Это важно, поскольку бесконтрольные повторные запуски могут повышать стоимость без реального улучшения результата.

\subsection{Связь прототипа с дальнейшей production-архитектурой}

Экспериментальный прототип должен быть спроектирован так, чтобы его результаты можно было использовать при дальнейшем развитии production-версии ClawdLab. Поэтому предлагаемые изменения должны быть совместимы с текущей архитектурной философией системы: серверный контроль над состоянием, явные контракты API, внешние агентные рантаймы и сохранение проверяемых артефактов.

Если эксперимент покажет положительный эффект, сущность TaskAttempt может стать постоянной частью модели данных ClawdLab. Она позволит фиксировать историю выполнения задачи, связывать повторные попытки с причинами пересмотра и строить более прозрачную траекторию исследования. Это полезно не только для экспериментов, но и для реального использования системы: пользователь сможет видеть не только финальный результат, но и путь, которым система к нему пришла.

Структурированная память также может быть постепенно развита в production-архитектуре. На первом этапе достаточно хранить явные связи между существующими объектами: задачами, попытками, источниками, артефактами, критикой и документами. На следующем этапе можно добавить правила переноса информации между рабочим, эпизодическим и каноническим уровнями. В дальнейшем поверх этого слоя возможно построить mind-map представление или более сложную графовую память.

Для production-сценариев также потребуется расширить наблюдаемость системы. В эксперименте фиксируются логи, попытки, артефакты и оценки. В промышленной версии эти данные могут использоваться для мониторинга качества, анализа ошибок, поиска регрессий и поддержки human-in-the-loop валидации. Это особенно важно для научных агентных систем, где доверие к результату зависит не только от финального ответа, но и от возможности восстановить процесс его получения.

Таким образом, ожидаемый Proof-of-Concept выполняет двойную функцию. С одной стороны, он служит экспериментальным инструментом для проверки гипотезы о влиянии памяти и цикличности на качество работы ClawdLab. С другой стороны, он задает минимальный архитектурный путь развития системы: от хранения отдельных артефактов к управляемому исследовательскому состоянию, в котором задачи, попытки, критика и выводы связаны между собой.

Источники:

1. Toward Ultra-Long-Horizon Agentic Science: Cognitive Accumulation for Machine Learning Engineering - [https://arxiv.org/abs/2601.10402](https://arxiv.org/abs/2601.10402)
2. Graph-Native Cognitive Memory for AI Agents - [https://arxiv.org/abs/2603.17244](https://arxiv.org/abs/2603.17244)
3. MindMap: Knowledge Graph Prompting Sparks Graph of Thoughts in Multi-Hop Reasoning - [https://arxiv.org/abs/2308.09729](https://arxiv.org/abs/2308.09729)
4. Towards an AI Co-Scientist - [https://arxiv.org/abs/2502.18864](https://arxiv.org/abs/2502.18864)
5. Agent Laboratory: Using LLM Agents as Research Assistants - [https://arxiv.org/abs/2501.04227](https://arxiv.org/abs/2501.04227)
